\documentclass[10pt,a4paper]{amsart}
\usepackage[margin=19mm]{geometry}
\usepackage{amsmath,amssymb,mathtools}
\usepackage[scr=rsfs,cal=euler]{mathalfa}
\usepackage{xfrac}
\usepackage[unicode,hidelinks,bookmarks=false]{hyperref}
\newcommand{\ie}{i.e.\ }
\newcommand{\cf}{cf.\ }
\newcommand{\resp}{respectively\ }
\newcommand{\Subsection}{\subsection}
\providecommand{\nobreakdash}{\nobreak\hbox{-}}
\setlength{\emergencystretch}{2em}
\multlinegap0pt
\usepackage[shortlabels]{enumitem}
\usepackage{amssymb}
\usepackage{tikz}
\usepackage{mathtools}
\newtheorem{lem}{Lemma}
\usepackage{cleveref}
\crefname{lem}{Lemma}{Lemmas}
\DeclareMathOperator{\Mon}{Mon}
\newcommand{\Z}{\mathbb{Z}}
\DeclareMathOperator{\expsum}{expsum}
\title{Undecidable Diophantine problems in \\ generalisations of one-relator groups}
\author{Robert D. Gray, Alex Levine}
\date{Source: arXiv:2605.30535v1 (CC BY 4.0)}
\begin{document}
\maketitle

\noindent\emph{Source and licence.} Undecidable Diophantine problems in \\ generalisations of one-relator groups. Version arXiv:2605.30535v1 (CC BY 4.0), distributed by arXiv under the Creative Commons Attribution 4.0 International licence, http://creativecommons.org/licenses/by/4.0/, as recorded in the arXiv licence metadata for this exact version and shown on the abstract page https://arxiv.org/abs/2605.30535v1. Full licence notice: \texttt{licenses/arxiv-2605.30535-CC-BY-4.0.txt}.\par\smallskip

\noindent\textit{Editorial scope.} Counted unit: the article's lem GBS-defable (source0.tex lines 629--658). The article's complete original proof of it is delivered with it (source0.tex lines 660--793, one \texttt{proof} environment of 134 lines). No mathematics is written, repaired or re-derived: the statement and the proof are the article's own lines, in the article's own notation, and no retained line is substituted or re-flowed.  Retained uncounted as the source-local context the proof needs: lines 367--391; lines 530--541.  Nothing was omitted from any retained span, so the package carries no omission record.  No retained line contains a citation, so the wrapper carries no bibliography and no bibliography entry.  No retained line contains a figure environment, an image-inclusion command or drawing code, so no image is delivered and the package declares no asset dependency.  Editorial formatting only, in addition: an AMS \texttt{amsart} preamble that restates the article's own declarations and macros used by the retained lines, namely the environments \texttt{lem} on separate counters, together with the standard packages those lines need.  The retained environments print in this document as Lemma 1 in order of appearance, whereas the article prints the same items with the numbering of its own full document.  The complete native source of the article as distributed in the arXiv e-print for this exact version is bundled unchanged as \texttt{sources/201678/source0.tex}.
	\begin{lem}
		\label{GBS-norm-form}
		Let \(G = \langle a, s, t \mid sas^{-1} = a^p, tat^{-1} = a^q
		\rangle\), where \(p\) and \(q\) are distinct primes. 
Define
				the following sets of freely reduced words:
				\begin{align*} 
T_s & =    
\{ s^{-i} a^j s^k \mid i,k \in \Z_{> 0}, \; j \in \{0, \ldots, p^i - 1\}, \; j \not\in p\Z \}  \\
& \cup \{ s^k \mid k \in \Z_{\geq 0} \} \\ 
& \cup \{ s^{-i} a^j \mid i \in \Z_{\geq 0}, \; j \in \{0, \ldots, p^i - 1\} \}, \\ 
T_t & =    
\{ t^{-i} a^j t^k \mid i,k \in \Z_{> 0}, \; j \in \{0, \ldots, q^i - 1\}, \; j \not\in q\Z \}  \\
& \cup \{ t^k \mid k \in \Z_{\geq 0} \} \\ 
& \cup \{ t^{-i} a^j \mid i \in \Z_{\geq 0}, \; j \in \{0, \ldots, q^i - 1\} \}.   
\end{align*} 

		Then 
		\[
			\{a^r x_1 y_1 \cdots x_n y_n \mid r \in \Z, n \in \Z_{\geq 0},
					x_1 \in T_s, x_2, \ldots, x_n \in T_s \setminus \{1\}, 
			y_1, \ldots, y_{n - 1} \in T_t \setminus \{1\}, y_n \in T_t \}
		\]
		is a normal form for \(G\).
	\end{lem}

\begin{lem}
		\label{GBS-cents}
		Let \(p, q \in \mathbb{Z}\) be distinct prime numbers. Define 
		\(G = \langle a, s, t \mid sas^{-1} = a^p, tat^{-1} = a^q \rangle\).
		Then the following hold in \(G\):
		\begin{enumerate}
			\item \(C_G(t) = \langle t \rangle\) and \(C_G(s) = \langle s \rangle\);
					\label{GBS-cent-s-t}
			\item For all \(i \in \Z_{\geq 0}\), \(C_G(st^i) = \langle st^i \rangle\).
					\label{GBS-cent-sti}
		\end{enumerate}
	\end{lem}

	\begin{lem}
		\label{GBS-defable}
		Let \(p, q \in \mathbb{Z}\) be distinct prime numbers.
		The following sets are equationally definable in
		\(G = \langle a, s, t \mid sas^{-1} = a^p, tat^{-1} = a^q \rangle\):
		\begin{enumerate}
				\item \(\Mon\langle s \rangle\) and \(\Mon\langle t \rangle\);
						\label{GBS-defable-monst}
				\item \(\{(g, t^i) \mid  i \in \Z_{\geq 0}, g \in \Mon\langle st^i
						\rangle\}\); 
						\label{GBS-defable-g-ti}
		\end{enumerate}
		In addition, if \(E\) is equationally definable
				in \(G\) and admits two coordinates that only take values in
				\(\Mon\langle s, t \rangle\), then the following three sets
				are equationally definable in \(G\), where the intersections
				are taken between the two coordinates in \(E\) taking
				values in \(\Mon \langle s, t \rangle\) and the given sets:
		\begin{enumerate}[resume]
				\item \(E \cap \{(u_1, u_2) \mid u_1, u_2 \in \Mon\langle s, t \rangle,
						\expsum_{s}(u_1) = \expsum_t(u_2)\}\); 
						\label{GBS-defable-expsum-s-t}
				\item \(E \cap \{(u_1, u_2) \mid u_1, u_2 \in \Mon\langle s, t \rangle,
						\expsum_{s}(u_1) = \expsum_s(u_2)\}\); 
						\label{GBS-defable-expsum-s-s}
				\item \(E \cap \{(u_1, u_2) \mid u_1, u_2 \in \Mon\langle s, t \rangle,
				\expsum_{t}(u_1) = \expsum_t(u_2)\}\). 
						\label{GBS-defable-expsum-t-t}
		\end{enumerate}
	\end{lem}

	\begin{proof}
			For \eqref{GBS-defable-monst} and \eqref{GBS-defable-g-ti}, using \cref{GBS-cents}, we will begin by showing
			that the set of solutions to
			\begin{equation}
					\label{eqn:GBS-Monsti}
					(X st^i = st^i X) \wedge (XaX^{-1} t^{-(i + 1)} at^{i + 1} = t^{-(i + 1)}
			at^{i + 1} XaX^{-1})
			\end{equation}
			is precisely \(\Mon\langle st^i \rangle\), for all \(i \in \Z_{\geq
			0}\). The fact that \(\Mon \langle s \rangle\) is equationally
			definable is the special case of this when \(i = 0\).

			Fix \(i \geq 0\), and let \(g \in 
			 C_G(st^i)\); that is, \(g\) is a solution to the first above equation.
			 By \cref{GBS-cents}(2), \(C_G(st^i) = \langle st^i \rangle\), and so
			 \(g = (st^i)^j\) for some \(j \in \Z\). We consider the cases when
			 \(j \geq 0\) and \(j < 0\) separately.

			 If \(j \geq 0\) (that is, \(g \in \Mon\langle st^i \rangle\)),
			 then \(gag^{-1} \in \langle a \rangle\). Further, \(t^{-(i + 1)} a
			 t^{i + 1} \cdot a = t^{-(i + 1)} a \cdot a^{p^{i + 1}} t^{i + 1} =
			 a \cdot t^{-(i + 1)} at^{i + 1}\), and so \(gag^{-1} t^{-(i + 1)}
			 a t^{i + 1} = t^{-(i + 1)} a t^{i + 1} gag^{-1}\). Thus
			 \((st^i)^j\) is a solution to the system Equation~\eqref{eqn:GBS-Monsti}
			 for all \(j \geq 0\).

	Next suppose that \(j <
			0\).  Then
			\begin{align*}
					gag^{-1} t^{-(i + 1)} a t^{i + 1} & = (st^i)^j a (st^i)^{-j} t^{-(i + 1)} a t^{i + 1} \\
										& = (st^i)^{j + 1} t^{-i} s^{-1} a s
										t^i (st^i)^{-j - 1} t^{-(i + 1)} a t^{i + 1}
										\\
										& = 
										\begin{cases}
										t^{-i} s^{-1} a s
										 t^{-1} a t^{i + 1} & j = -1 \\
										(st^i)^{j + 1} t^{-i} s^{-1} a s
										t^i (st^i)^{-j -2} s t^{-1} a t^{i + 1}
										& j < -1,
										\end{cases}
										\\
										t^{-(i + 1)} a t^{i + 1} gag^{-1}
										& = t^{-(i + 1)} a t^{i + 1}
										(st^i)^j a (st^i)^{-j} \\
										& = t^{-(i + 1)} a t^{i + 1}
										(st^i)^{j + 1} t^{-i} s^{-1} a s
								t^i (st^i)^{-j - 1} \\
										& = \begin{cases}
										t^{-(i + 1)} a t s^{-1} a s t^i & j = -1 \\
										t^{-(i + 1)} a t s^{-1}
										(st^i)^{j + 2} t^{-i} s^{-1} a s
								t^i (st^i)^{-j - 1} 
																		& j < -1.
										\end{cases}
			\end{align*}
			We first consider the case that \(j = -1\). We can see that
			\(t^{-i}, t^{-1} a t^{i + 1} \in T_t \setminus \{1\}\) and \(s^{-1}
			a s\in T_s \setminus \{1\}\) (using the definitions of \(T_s\) and
			\(T_t\) from \cref{GBS-norm-form}) and so \(t^{-i} s^{-1} a s
			t^{-1} a t^{i + 1}\) is in normal form. Similarly we have that
			\(t^{-(i + 1)} a t, \; t^i \in T_t \setminus \{1\}\) and
			\(s^{-1} a s \in T_s \setminus \{1\}\), and so
			\(t^{-(i + 1)} a t s^{-1} a st^i\) is also in normal form.
			Thus the group elements these words represent coincide
			if and only if these words are equal as words, which they do not,
			and so this case of \(j = -1\) contributes no solutions to
			the system Equation~\eqref{eqn:GBS-Monsti}.

			Similarly, if \(j < -1\), then \(s, s^{-1} as \in T_s \setminus
			\{1\}\) and \(t^i, t^{-i}, t^{-1} a t^{i + 1} \in T_t \setminus
			\{1\}\), and so the word \((st^i)^{j + 1} t^{-i} s^{-1} a s t^i
			(st^i)^{-j -2} s t^{-1} a t^{i + 1}\) is in the normal form from
			\cref{GBS-cents}. One can see similarly that \(t^{-(i + 1)} a t
			s^{-1} (st^i)^{j + 2} t^{-i} s^{-1} a s t^i (st^i)^{-j - 1}\) is in
			normal form, and so these elements coincide if and only if these
			words are equal. It is immediate that this never happens when \(j <
			0\), and so this case also does not contribute solutions to the
			system Equation~\eqref{eqn:GBS-Monsti}.

We have thus shown that the set of solutions to the equation
			Equation~\eqref{eqn:GBS-Monsti} is \(\{(st^i)^j \mid j \in \Z_{\geq 0}\} =
			\Mon \langle st^i \rangle\), and so \(\Mon\langle st^i \rangle\) is
			equationally definable in \(G\) for all \(i \in \Z_{\geq 0}\). A
			particular case of this is when \(i = 0\), which shows that
			\(\Mon \langle s \rangle\) is equationally definable. By
			symmetry, we obtain that \(\Mon \langle t \rangle\) is also
			equationally definable and we have shown \eqref{GBS-defable-monst}.

			To obtain the more
			general desired set for \eqref{GBS-defable-g-ti}, note that if \(\mathcal E\) is a
			system of equations in \(G\) with a variable \(Y\) admitting the
			set of solutions \(\Mon\langle t \rangle\), then we have shown
			above that the set of solutions to \(\mathcal E \wedge (XaX^{-1}
			t^{-1} Y^{-1} a Y t = t^{-1} Y^{-1} a Y t XaX^{-1}) \wedge (sY X =
			XsY)\) is precisely the desired set stated in \eqref{GBS-defable-g-ti}.

For \eqref{GBS-defable-expsum-s-s} (\eqref{GBS-defable-expsum-t-t}
			is symmetric), consider the set of solutions to \((X, Y)\) in \( X
			U a(XU)^{-1} = Y a (Y)^{-1}\), with the constraints that \((X, Y)
			\in E\) and \(U \in \Mon \langle t \rangle\).  So let \((x, y) \in
			E \subseteq (\Mon \langle s, t \rangle)^2\), and \(i \in \Z_{\geq
			0}\). Then \(xt^i a (xt^i)^{-1} = a^{p^{\expsum_s(x)}q^{i +
			\expsum_t(x)}}\) and \(y a y^{-1} = a^{p^{\expsum_s(y)}
			q^{\expsum_t(y)}}\). Since \(p\) and \(q\) are distinct primes,
			this holds if and only if \(\expsum_s(x) = \expsum_s(y)\) and \(i +
			\expsum_t(s) = \expsum_t(y)\).  Thus the set of \(x\), \(y\) and
			\(i\) for which this is satisfied is
			\[
				\{(x, y, t^i) \in E \times \Mon\langle t^i \rangle \mid
				\expsum_s(x) = \expsum_s(y) \text{ and } \expsum_t(x) =
		\expsum_t(y) + i\},
			\]
			and so this is the set of solutions to \((X, Y, U)\) in this system.
			Thus if we simply consider the solutions to \((X, y)\) in this
			system, we obtain the projection of the above set to the first two
			coordinates, which is:
			\[
				\{(x, y) \in E \times \Mon\langle t^i \rangle \mid
		\expsum_s(x) = \expsum_s(y)\},
			\]
	which is the desired set.
			
	For \eqref{GBS-defable-expsum-s-t}, let \(\mathcal E\) be a system of equations in \(G\) with
			a variable \(U\), such that the set of solutions to \(U\) is \(\Mon
			\langle st \rangle\), which exists by \eqref{GBS-defable-g-ti}.  Let \(X\) and \(Y\) be
			additional variables with the constraints that \(X, Y \in \Mon
			\langle s, t \rangle\). We can then use \eqref{GBS-defable-expsum-s-s} and \eqref{GBS-defable-expsum-t-t} to insist
			that \(\expsum_s(X) = \expsum_s(U)\) and \(\expsum_t(Y) = 
			\expsum_t(U)\). Since \(U \in \Mon\langle st \rangle\),
			\(\expsum_s(U) = \expsum_t(U)\), and so this system of
			equations will have the desired set of solutions to
			\((X, Y)\).
	\end{proof}

\end{document}
